\documentclass[11pt]{article}

% ===================== PACKAGES =====================
\usepackage[a4paper,margin=1in]{geometry}
\usepackage{amsmath,amssymb}
\usepackage{enumitem}
\usepackage{fancyhdr}
\usepackage{xcolor} 

% ===================== HEADER & FOOTER =====================
\pagestyle{fancy}
\fancyhf{}
\lhead{CSE 400: Fundamentals of Probability in Computing}
\rhead{Lecture-3 Scribe Submission}
\cfoot{\thepage}

% ===================== CUSTOM COMMANDS =====================
\newcommand{\solution}{
    \vspace{0.5cm}
    \noindent\textbf{\textcolor{blue}{Solution:}}
    \vspace{0.2cm}
    
}

% ===================== TITLE =====================
\title{
    \normalsize School of Engineering and Applied Science (SEAS), Ahmedabad University \\
    \vspace{0.2cm}
    \textbf{CSE 400: Fundamentals of Probability in Computing}\\
    \Large Lecture-3 Scribe Submission
}
\author{} 
\date{}

\begin{document}
\maketitle

\vspace{-2cm}
\begin{center}
    \begin{tabular}{ll}
        \textbf{Name:} & {Jashvi Shah \hspace{2.5in}} \\ [1.5ex]
        \textbf{Enrollment No:} & {AU2440026 \hspace{2.5in}} \\ [1.5ex]
        \textbf{Email:} & {jashvi.s@ahduni.edu.in \hspace{2.1in}} \\ [1.5ex]
        \textbf{Date of Submission:} & {Febrauary 7\textsuperscript{th}, 2026 (11:59 PM) \hspace{1.2in}}
    \end{tabular}
\end{center}

\hrule
\vspace{0.5cm}

\section*{CSE400 – Lecture 3 Scribe}
\subsection*{Project Kickoff and Introduction to Undergraduate Research Programme (UGRP)}

\section*{1. Objective of the Lecture}

Lecture 3 introduces:

\begin{itemize}
\item The CSE400 project framework (30\%)
\item Team formation and execution workflow
\item Milestone-based project development (M1–M6)
\item Scribe and artifact expectations
\item Introduction to UGRP (Undergraduate Research Programme)
\item MICxN Research Lab ecosystem and outcomes
\end{itemize}

The lecture establishes how students transition from coursework to research-driven, milestone-oriented project execution.

\section*{2. CSE400 Project Structure}

\subsection*{Weightage}

The semester project carries 30\% of the total course evaluation.

\subsection*{Team Formation}

Teams must be formed by

17 January 2026 (Saturday – End of Day).

This deadline marks the start of formal project execution.

\section*{3. Project Execution Guidelines}

Each group must complete six milestone submissions (M1–M6) during the semester.

Additionally, multiple artifact types are required:

\begin{itemize}
\item Concept Evolution Maps
\item Scribe: Process \& Decision Making
\item Multimodal Artifacts (Video / Audio / Visual)
\item Question-Driven Artifacts
\item Collaboration \& Team Dynamics Artifacts
\end{itemize}

Deliverables include:

\begin{itemize}
\item Codes
\item Reports
\item Videos
\item Other materials as specified.
\end{itemize}

Evaluation includes:

\begin{itemize}
\item Team assessment (before and after mid-semester)
\item Final project viva
\item Final submission at semester end.
\end{itemize}

\section*{4. Step-by-Step Project Workflow (M1–M6)}

The lecture defines a sequential execution pipeline.

\subsection*{M1 – Kickstart}

Includes:

\begin{itemize}
\item Team formation
\item Area identification
\item Background study
\item Motivation
\item Problem formulation
\end{itemize}

This stage converts a general interest area into a clearly defined problem.

\subsection*{M2 – Mathematical Modeling}

Students mathematically model their chosen domain using:

\begin{itemize}
\item Random Variables (RV)
\item PMF / PDF
\item CDF
\item Multivariate Random Variables
\item Joint PMF / PDF / CDF
\end{itemize}

Purpose: convert real-world problems into probabilistic representations.

\subsection*{M3 – Coding}

Implementation phase involving:

\begin{itemize}
\item Simulation
\item Computation
\end{itemize}

Models developed in M2 are translated into executable code.

\subsection*{M4 – Inference}

Students must:

\begin{itemize}
\item Choose a randomized algorithm
\item Understand its mechanism
\item Implement it
\end{itemize}

This milestone introduces algorithmic reasoning.

\subsection*{M5 – Randomized Algorithms}

The selected randomized algorithm is applied to the domain problem.

Students compare:

Randomized results vs Deterministic approaches

Goal: derive new inferences.

\subsection*{M6 – Final Analysis}

Includes:

\begin{itemize}
\item Deriving bounds
\item Performing analysis
\item Compiling final deliverables
\end{itemize}

This completes the research pipeline.

\section*{5. Submission Types}

\subsection*{Submission \#1 – Concept Evolution Maps}

Used to visually organize ideas.

Tools:

\begin{itemize}
\item Miro
\item draw.io
\end{itemize}

\subsection*{Submission \#2 – Scribe (Learning Reflection Logs)}

Two types:

\begin{itemize}
\item Lecture Scribe
\item Project Scribe
\end{itemize}

Lecture scribes:

\begin{itemize}
\item Assigned to two groups per lecture
\item Must reflect lecture content
\item Minimum 8–10 pages
\end{itemize}

Project scribes must document:

\begin{itemize}
\item Decision logs (Why X over Y)
\item Constraints
\item Alternatives considered
\item Final decisions
\item Evidence used
\item Trade-off matrices (Cost vs Performance vs Risk)
\end{itemize}

Total: 6 scribe submissions (bi-weekly).

\subsection*{Submission \#3 – Multimodal Artifacts}

Includes:

\begin{itemize}
\item Think-aloud videos
\item One-minute insight videos
\item Project demos
\end{itemize}

Rules:

\begin{itemize}
\item One video per milestone
\item Duration: $\sim$10–15 minutes
\item Content matters more than editing
\item PPT/Google Slides may be used
\end{itemize}

\section*{6. Undergraduate Research Programme (UGRP)}

UGRP promotes the idea of a T-shaped individual:

Vertical bar $\rightarrow$ depth in one discipline

Horizontal bar $\rightarrow$ ability to work across disciplines

This model emphasizes both specialization and collaboration.

\section*{7. Philosophy of UGRP}

UGRP is based on:

\begin{itemize}
\item Multidisciplinary learning (Arts, Science, Management)
\item Experiential learning
\item Research-driven education
\end{itemize}

It follows a 4D model:

\begin{itemize}
\item Discover
\item Design
\item Develop
\item Deliver
\end{itemize}

Focus areas:

\begin{itemize}
\item CS and CSE
\item Data Science and Applied AI
\item Modern Computer Systems (Hardware + Software)
\item Networks and IoT / IoBNT / IoV
\end{itemize}

\section*{8. MICxN Research Lab Activities}

Key industry projects include:

\begin{itemize}
\item 5G-enabled Intelligent Transportation Systems (ITS) Testbed in Gujarat
\item AI-based ITS solutions
\item V2X communication between OBU and RSU
\item SDK testing
\item Sensor data collection
\item In-car dashboard development
\end{itemize}

Research directions include:

\begin{itemize}
\item Deep Learning crash prevention
\item Deep-RL smart signaling
\item mmWave beam prediction
\item C-V2X and Wi-Fi coexistence
\item Testbed development is conducted at Ahmedabad University.
\end{itemize}

\section*{9. Research Collaborations and Outcomes}

MICxN collaborates with international universities including:

\begin{itemize}
\item University of Liverpool
\item University of Manchester
\item NTU Singapore
\item Boston College
\item Monash University
\item Qatar University
\item University of Oulu
\end{itemize}

UGRP students have:

\begin{itemize}
\item Published in IEEE venues
\item Received Best Paper and Best UG Presentation awards
\item Joined Amazon, Google, Intel, American Express
\item Pursued PhDs internationally
\end{itemize}

Multiple papers were accepted at COMSNETS and NCC conferences.

\section*{10. Current MICxN Lab Scholars}

Students work on:

\begin{itemize}
\item Advanced spectrum sensing
\item O-RAN
\item 6G ISAC + UAV
\item Wi-Fi sensing
\item Wireless networks
\item Wi-Fi 8
\item Antenna design
\end{itemize}

\section*{11. Conclusion}

Lecture 3 establishes:

\begin{itemize}
\item A milestone-driven project framework
\item Artifact-based assessment
\item Integration of research into undergraduate learning
\item Real-world impact through MICxN lab and UGRP
\end{itemize}

The session concludes with an open Q\&A.

\end{document}
